\documentclass[11pt,a4paper]{ctexart}
\usepackage[a4paper,margin=1.78cm,headheight=15pt]{geometry}
\usepackage{amsmath,amssymb,mathtools,bm}
\usepackage{booktabs,longtable,tabularx,array}
\usepackage{enumitem}
\usepackage[most]{tcolorbox}
\usepackage{tikz}
\usetikzlibrary{arrows.meta,positioning,fit,calc}
\usepackage{xcolor}
\usepackage{fancyhdr}
\usepackage{hyperref}
\usepackage{microtype}

\newcolumntype{Y}{>{\raggedright\arraybackslash}X}
\newcolumntype{L}[1]{>{\raggedright\arraybackslash}p{#1}}
\setlength{\parindent}{2em}
\setlength{\parskip}{0.34em}
\setlength{\emergencystretch}{3em}
\renewcommand{\arraystretch}{1.22}
\setlength{\tabcolsep}{3pt}
\setlist[itemize]{itemsep=0.18em,topsep=0.35em,leftmargin=2em}
\setlist[enumerate]{itemsep=0.18em,topsep=0.35em,leftmargin=2em}

\definecolor{deep}{RGB}{42,58,73}
\definecolor{soft}{RGB}{245,247,249}
\definecolor{line}{RGB}{145,154,163}
\definecolor{accent}{RGB}{104,76,55}
\definecolor{greenish}{RGB}{57,92,76}
\definecolor{warn}{RGB}{136,68,63}
\definecolor{purple}{RGB}{87,72,116}

% ── 双主题图解语义色彩（LaTeX 打印高对比白底）──
\definecolor{themebg}{HTML}{FFFFFF}
\definecolor{themefg}{HTML}{111111}
\definecolor{themegray}{HTML}{666666}
\definecolor{themecurve}{HTML}{1D63B8}
\definecolor{themealert}{HTML}{C53030}
\definecolor{themeamber}{HTML}{B86E00}
\definecolor{themeblue}{HTML}{1D63B8}

\newenvironment{handbookdiagram}{\par\vspace{0.4em}\noindent\begin{center}}{\end{center}\vspace{0.4em}}
\newcommand{\diagramcaption}[1]{\par\vspace{0.3em}{\small\textbf{#1}}}

\pgfdeclarelayer{background}
\pgfdeclarelayer{foreground}
\pgfsetlayers{background,main,foreground}

\hypersetup{
  colorlinks=true,
  linkcolor=deep,
  urlcolor=accent,
  pdftitle={数字特征与依赖摘要}
}

\pagestyle{fancy}
\fancyhf{}
\lhead{\small\color{deep}\textbf{数学一 · 概率统计 Topic 05}}
\rhead{\small\color{deep}期望 · 方差 · 协方差 · 相关系数 · 条件期望}
\cfoot{\small\thepage}
\renewcommand{\headrulewidth}{0.4pt}

\newtcolorbox{corebox}[1]{breakable,colback=soft,colframe=deep,title=\textbf{#1},boxrule=0.8pt,arc=2mm}
\newtcolorbox{methodbox}[1]{breakable,colback=white,colframe=greenish,title=\textbf{#1},boxrule=0.7pt,arc=1.5mm}
\newtcolorbox{warnbox}[1]{breakable,colback=white,colframe=warn,title=\textbf{#1},boxrule=0.7pt,arc=1.5mm}
\newtcolorbox{examplebox}[1]{breakable,colback=white,colframe=purple,title=\textbf{#1},boxrule=0.7pt,arc=1.5mm}
\newcommand{\kw}[1]{\textbf{\color{deep}#1}}
\newcommand{\indep}{\mathrel{\perp\!\!\!\perp}}

\title{\vspace{-1.1cm}\textbf{数字特征与依赖摘要}\\[0.4em]
\large 怎样高效压缩概率分布而不遗忘关键结构与信息损失？}
\author{数学一 · 概率统计 Topic 05}
\date{}

\begin{document}
\maketitle

\begin{corebox}{母问题：怎样用少量数字摘要分布，同时始终知道自己丢掉了什么信息？}
完整分布通常包含远多于少数几个矩的信息。期望、方差、协方差和相关系数不是脱离任务意义上的“最佳压缩”，而是针对\textbf{位置、波动与线性依赖}的有损摘要：它们便于计算，却不能一般恢复分布或完整依赖结构。

因此本册不是一条“$E\to D\to\operatorname{Cov}\to\rho$”的生成流水线，而是两类摘要从各自概率对象分叉：
\[
\boxed{
\begin{array}{c}
P_X\ \longrightarrow\ \{E(X),\operatorname{Var}(X),\text{高阶矩/绝对矩}\}\\[0.5em]
P_{X,Y}\ \longrightarrow\ \operatorname{Cov}(X,Y)\ \longrightarrow\ \rho_{XY}
\end{array}}
\]
条件期望与全方差进一步回答：当信息分层以后，总体摘要怎样由“层内波动”和“层间差异”共同组成。
\end{corebox}

\section{本册定位与职责划分}

\begin{itemize}
  \item \kw{本册拥有（Owns）：}数学期望 $E(X)$、方差 $\operatorname{Var}(X)$、协方差 $\operatorname{Cov}(X,Y)$、相关系数 $\rho_{XY}$、矩（原点矩与中心矩）、条件期望，以及全期望、全方差与全协方差分解。
  \item \kw{本册使用（Uses）：}一维分布与二维联合分布。
  \item \kw{本册不拥有（Does not own）：}大数定律与中心极限定理的收敛证明（Topic 06 负责）、统计量样本方差的抽样分布（Topic 07 负责）。
  \item \kw{输出目标（Output）：}分布标量特征的完整算子性质与计算链条。
\end{itemize}

\section{第一层：期望、矩、形状与存在性}

数字特征是有损压缩：同一组均值和方差可以对应完全不同的分布。因此每次使用矩时，必须同时问“这个矩是否存在”和“它能表达什么信息”。

离散型与具有密度的连续型随机变量分别用
\[
E(X)=\sum_kx_kp_k,\qquad E(X)=\int_{-\infty}^{\infty}xf_X(x)\,dx,
\]
并要求 $E\lvert X\rvert<\infty$，才能把期望作为有限实数使用。更一般地，函数期望可以直接对原随机变量的分布求和/积分；本册把这条路线简称为 LOTUS（即“无需先求 $g(X)$ 的分布，直接计算 $E[g(X)]$”）：
\[
\boxed{E[g(X)]=\sum_kg(x_k)p_k\quad\text{或}\quad E[g(X)]=\int g(x)f_X(x)\,dx},
\]
它只调用 $X$ 的分布，不要求先求 $g(X)$ 的完整分布。期望是线性算子：只要各项存在，
\[
E\!\left(\sum_i a_iX_i+b\right)=\sum_i a_iE(X_i)+b,
\]
这里不需要独立性；独立性只在乘积拆分、卷积或协方差消失等更强动作中出现。

\begin{methodbox}{LOTUS 与线性的证据骨架}
离散情形把 $g(X)$ 取值相同的样本点合并：
\[
E[g(X)]=\sum_y yP(g(X)=y)=\sum_xg(x)P(X=x).
\]
连续情形把求和换成对 $X$ 的分布积分；因此只要目标是期望，就不必先构造 $g(X)$ 的完整分布。线性性来自积分/求和的可加性和常数可提出，整个推导没有使用 $X,Y$ 独立。使用前仍需检查相应绝对可积性，不能用形式线性掩盖发散。
\end{methodbox}
\begin{methodbox}{先找逐点恒等式，再决定是否需要独立性}
数字特征题的第一动作不是立即写分布，而是把目标随机变量在每个样本点上的表达式化简。例如
\[
\max(X,Y)\min(X,Y)=XY,
\qquad
\max(X,Y)=\frac{X+Y+\lvert X-Y\rvert}{2},
\qquad
\min(X,Y)=\frac{X+Y-\lvert X-Y\rvert}{2},
\]
以及
\[
(X-\mu)^2=X^2-2\mu X+\mu^2.
\]
前两条还给出
\[
\max(X,Y)+\min(X,Y)=X+Y,
\qquad
\max(X,Y)-\min(X,Y)=\lvert X-Y\rvert,
\]
因此求极值的期望时，往往只需处理 $E\lvert X-Y\rvert$，不必先构造极值的完整分布。
若题目只问 $E[\max(X,Y)\min(X,Y)]$，先用该恒等式和独立性得到 $E(XY)=E(X)E(Y)$；只有要求 $\max$、$\min$ 的分布或其他非线性输出时，才进入联合分布/顺序统计量路线。这样既避免无谓构造中间分布，也避免把“代数恒等式”误当成独立性结论。
\end{methodbox}

\begin{methodbox}{高阶矩先做归一化与分布识别：2010 Q十四}
若题目以级数形式给出离散分布，先检查总质量，再决定用逐项求和还是已知分布矩。题面
\[
P(X=k)=\frac{C}{k!},\qquad k=0,1,2,\ldots
\]
要求 $\sum_kP(X=k)=1$，故 $Ce=1$、$C=e^{-1}$，这正是 $X\sim\operatorname{Pois}(1)$。于是
\[
E(X^2)=\operatorname{Var}(X)+[E(X)]^2=1+1=2.
\]
也可直接用级数验证：$E(X^2)=e^{-1}\sum_{k\ge1}k^2/k!=e^{-1}(\sum k(k-1)/k!+\sum k/k!)=2$。第一种路径揭示生成机制，第二种路径可作为不依赖分布表的独立回检；不能在 $C$ 尚未归一化时直接套矩公式。
\end{methodbox}

\begin{examplebox}{只问数字特征时直接压缩：1990、1992、1995 真题}
若 $X\sim\operatorname{Pois}(2)$、$Z=3X-2$，线性性直接给出 $E(Z)=4$，不必先求 $Z$ 的完整分布。若 $X\sim\operatorname{Exp}(1)$，则
\[
E(X+e^{-2X})=E(X)+E(e^{-2X})=1+\int_0^\infty e^{-3x}\,dx=\frac43.
\]
若 $X\sim\operatorname{Bin}(10,0.4)$，二阶矩走方差分流：
\[
E(X^2)=\operatorname{Var}(X)+[EX]^2=10(0.4)(0.6)+4^2=18.4.
\]
三题的共同决策是先看目标是否只要求数字特征；若是，优先 LOTUS、线性性或二阶矩恒等式，停止构造不必要的中间分布。
\end{examplebox}

\begin{examplebox}{2011 真题：乘积高阶矩先拆独立，再辨认原点矩}
设 $X,Y$ 相互独立，且分别服从 $N(\mu,\sigma^2)$，求 $E(XY^2)$。这里目标是乘积矩，不是线性组合；第一步必须使用独立性：
\[
E(XY^2)=E(X)E(Y^2).
\]
正态变量的二阶\textbf{原点矩}为
\[
E(Y^2)=\operatorname{Var}(Y)+[E(Y)]^2=\sigma^2+\mu^2,
\qquad E(X)=\mu,
\]
所以
\[
\boxed{E(XY^2)=\mu(\mu^2+\sigma^2)}.
\]
易错边界是把 $E(Y^2)$ 误写成方差 $\sigma^2$；只有中心二阶矩 $E[(Y-\mu)^2]$ 才等于 $\sigma^2$。若 $X,Y$ 不独立，边缘正态分布本身不足以确定 $E(XY^2)$，必须回到联合分布。
\end{examplebox}

\begin{examplebox}{中心化成对和：2001 Q十二 的协方差展开}
令 $A_i=X_i+X_{n+i}$，全样本容量为 $2n$。则
\[
E(A_i-2\bar X)=0,\quad D(A_i)=2\sigma^2,\quad D(2\bar X)=2\sigma^2,
\quad \operatorname{Cov}(A_i,2\bar X)=\frac{2\sigma^2}{n}.
\]
因此
\[
D(A_i-2\bar X)=2\sigma^2\left(1-\frac1n\right),
\qquad
E\sum_{i=1}^n(A_i-2\bar X)^2=2(n-1)\sigma^2.
\]
这里目标只是期望，不需要求整组括号项的联合分布；但不能把它们当成独立项，因为它们共享 $\bar X$ 并受中心化约束。
\end{examplebox}

\subsection{原点矩、中心矩与绝对矩}

对 $r>0$，定义 $r$ 阶原点矩、绝对矩和中心矩：
\[
\mu'_r=E(X^r),\qquad m_r=E(\lvert X\rvert^r),\qquad
\mu_r=E[(X-\mu)^r],\quad \mu=E(X).
\]
二阶中心矩就是方差：$\mu_2=\operatorname{Var}(X)$；一阶中心矩恒为 $0$。标准化矩为
\[
\eta_r=E\left[\left(\frac{X-\mu}{\sigma}\right)^r\right],\qquad 0<\sigma^2<\infty.
\]
三阶标准化中心矩是偏度，四阶标准化中心矩是峰度：
\[
\gamma_1=\frac{\mu_3}{\sigma^3},\qquad
\beta_2=\frac{\mu_4}{\sigma^4},\qquad \gamma_2=\beta_2-3.
\]
正态分布的 $\gamma_1=0,\gamma_2=0$；$\gamma_1=0$ 不推出对称，峰度也不只是“峰顶高度”。

若 $0<r<s$ 且 $E\lvert X\rvert^s<\infty$，则由 Hölder/Lyapunov 不等式
\[
\left(E\lvert X\rvert^r\right)^{1/r}\leq\left(E\lvert X\rvert^s\right)^{1/s}.
\]
所以高阶绝对矩存在会推出低阶绝对矩存在，反向一般不成立。Cauchy 分布就是典型边界：位置参数存在，但一阶期望和二阶方差均不存在，不能机械代入“均值—方差”公式。

\begin{methodbox}{绝对偏差先按中心分层}
目标 $E\lvert X-c\rvert$ 不是标准差，也不能由 $E(X)$ 或 $D(X)$ 单独确定。先确认中心 $c$ 是题目给定的常数还是 $E(X)$，再按 $X<c$、$X=c$、$X>c$ 分层：
\[
E\lvert X-c\rvert=E[(c-X)I_{\{X<c\}}]+E[(X-c)I_{\{X>c\}}].
\]
连续变量用两侧积分；离散变量按每个取值的质量求和。若 $X$ 为整数型且 $c$ 为整数，$X=c$ 项贡献为零，但不能把它从概率质量守恒中“删掉”；它只是在绝对偏差和中系数为零。最后检查结果非负，并区分精确数字特征与由方差给出的上界。
\end{methodbox}

\begin{examplebox}{2023 真题：Poisson(1) 的绝对偏差}
若 $X\sim\operatorname{Pois}(1)$，则 $EX=1$。按 $0,1,\{2,3,\ldots\}$ 三层记账：
\[
\begin{aligned}
E\lvert X-1\rvert
&=P(X=0)+\sum_{k=2}^{\infty}(k-1)P(X=k)\\
&=e^{-1}+e^{-1}\sum_{k=2}^{\infty}\frac{k-1}{k!}
=2e^{-1}.
\end{aligned}
\]
也可用 $E\lvert X-1\rvert=E(X-1)+2E[(1-X)I_{\{X=0\}}]$ 回检同一结果。不能把该量写成 $\sqrt{D(X)}=1$；标准差与绝对偏差是不同的损失/摘要。
\end{examplebox}

\begin{examplebox}{2009 真题：二项总体中“均值减方差”构造 $np^2$}
若 $X_i\sim\operatorname{Bin}(n,p)$，$\bar X$ 与 $S^2$ 分别为样本均值和无偏样本方差，则
\[
E(\bar X)=np,\qquad E(S^2)=D(X)=np(1-p).
\]
目标参数可重写为
\[
np^2=np-np(1-p)=E(X)-D(X).
\]
因此
\[
E(\bar X-S^2)=np-np(1-p)=np^2,
\]
从而题目中的 $k=-1$。这不是偶然系数：$\bar X$ 提供总体均值摘要，$S^2$ 提供总体方差摘要；若把 $S^2$ 误当成 $np^2$ 的直接估计，就会丢失二项分布的方差项。该题只问无偏系数，在线性期望与 $E(S^2)=D(X)$ 闭合处停止，不需要构造估计量的完整抽样分布。
\end{examplebox}

\subsection{示性变量：把计数变成期望}

事件 $A$ 的示性变量 $I_A$ 满足
\[
I_A=\begin{cases}1,&A\text{发生},\\0,&A^c\text{发生},\end{cases}\qquad E(I_A)=P(A),\qquad I_A^2=I_A.
\]
若 $N=\sum_{i=1}^nI_{A_i}$ 是满足条件的对象数，则
\[
E(N)=\sum_{i=1}^nP(A_i)
\]
不需要事件独立。二阶矩则会引入交叉项：
\[
E(N^2)=\sum_iP(A_i)+2\sum_{i<j}P(A_i\cap A_j).
\]
这一区分解释了“期望易算而方差需处理依赖”的常见现象。

\subsection{尾和与尾积分}

若 $X\geq0$，Tonelli 定理给出
\[
\boxed{E(X)=\int_0^{\infty}P(X>t)\,dt}.
\]
若 $X$ 为非负整数值，
\[
\boxed{E(X)=\sum_{k=1}^{\infty}P(X\geq k)}.
\]
一般可积变量可拆为正负部分：
\[
E(X)=\int_0^\infty P(X>t)\,dt-\int_0^\infty P(X<-t)\,dt.
\]
等待时间、寿命、最大值和停止时刻问题，往往先写尾概率比直接求密度更短；但必须先确认非负性或绝对可积性。

\section{第二层：方差、标准化与最小二乘几何}

\subsection{方差的定义、性质与计算链}

当 $E(X^2)<\infty$ 时，
\[
\operatorname{Var}(X)=E[(X-\mu)^2]=E(X^2)-\mu^2\geq0,\qquad \operatorname{SD}(X)=\sqrt{\operatorname{Var}(X)}.
\]
这里先定义“几乎处处”：在概率语境下，性质对 $X$ 几乎处处成立，是指不满足该性质的样本点集合概率为 $0$；例如 $X=\mu$ 几乎处处等价于 $P(X=\mu)=1$。因此方差为零当且仅当 $X=\mu$ 几乎处处。平移不改变方差，缩放使方差乘以系数平方：
\[
\operatorname{Var}(aX+b)=a^2\operatorname{Var}(X).
\]
对有限个二阶可积变量，
\[
\operatorname{Var}\left(\sum_{i=1}^na_iX_i\right)
=\sum_{i=1}^na_i^2\operatorname{Var}(X_i)
+2\sum_{i<j}a_ia_j\operatorname{Cov}(X_i,X_j).
\]
只有在独立或至少两两不相关时，交叉项才能消失；这里“两两不相关”按定义指任意 $i\ne j$ 都有 $\operatorname{Cov}(X_i,X_j)=0$。

\begin{methodbox}{方差恒等式的计算推导}
令 $\mu=E(X)$，对任意常数 $a$ 展开平方：
\[
E[(X-a)^2]
=E[(X-\mu)+(\mu-a)]^2
=\operatorname{Var}(X)+(\mu-a)^2,
\]
因为 $E(X-\mu)=0$，交叉项恰好消失。取 $a=0$ 得
$\operatorname{Var}(X)=E(X^2)-\mu^2$；对和式逐项展开则得到协方差交叉项。这个推导同时说明了方差非负、均值的平方损失最优性和“方差不是线性算子”的边界。
\end{methodbox}

\subsection{均值是平方损失下的最优常数}

对任意常数 $a$，
\[
\boxed{E[(X-a)^2]=\operatorname{Var}(X)+(\mu-a)^2}.
\]
因此 $a=\mu$ 唯一（除退化情形外）最小化均方误差；这不是说均值必然是最可能取值，而是说它对应平方损失的最佳常数预测。

若 $0<\sigma^2<\infty$，标准化 $Z=(X-\mu)/\sigma$ 满足 $E(Z)=0,\operatorname{Var}(Z)=1$，但不会因此变成标准正态。变异系数
\[
\operatorname{CV}(X)=\frac{\sigma}{\lvert\mu\rvert}\quad(\mu\neq0)
\]
是无量纲的相对波动指标；均值接近零或允许负均值时，解释必须谨慎。

\section{第三层：协方差、相关与矩阵}

\subsection{协方差的双线性结构}

若二阶矩存在，
\[
\operatorname{Cov}(X,Y)=E[(X-EX)(Y-EY)]=E(XY)-E(X)E(Y).
\]
它对称、双线性，且 $\operatorname{Cov}(X,X)=\operatorname{Var}(X)$、$\operatorname{Cov}(X,c)=0$。特别地，
\[
\operatorname{Var}(X\pm Y)=\operatorname{Var}(X)+\operatorname{Var}(Y)\pm2\operatorname{Cov}(X,Y).
\]

由 Cauchy--Schwarz，
\[
\lvert\operatorname{Cov}(X,Y)\rvert\leq\sigma_X\sigma_Y.
\]
当两边方差正且有限时，相关系数
\[
\rho_{XY}=\frac{\operatorname{Cov}(X,Y)}{\sigma_X\sigma_Y}\in[-1,1].
\]
$\lvert\rho\rvert=1$ 当且仅当 $Y=aX+b$ 几乎处处成立，其中 $a$ 的正负决定相关系数的符号。

\subsection{独立、不相关与非线性依赖}

独立且二阶矩存在 $\Rightarrow \operatorname{Cov}(X,Y)=0$，但逆命题一般不成立。例：$X\sim U(-1,1),Y=X^2$ 时，奇偶对称使 $E(XY)=E(X^3)=0$，故不相关；然而 $Y$ 完全由 $X$ 决定，显然不独立。只有联合正态这一特殊结构把“不相关”升级为“独立”。

\begin{handbookdiagram}
\resizebox{0.92\linewidth}{!}{%
% ── V12 零协方差与非独立反例（大空间呼吸感 2 行布局） ──
\begin{tikzpicture}[
  scale=1.0,
  >=Stealth,
  font=\small,
  color=themefg,
  text=themefg
]

\pgfdeclarelayer{background}
\pgfdeclarelayer{foreground}
\pgfsetlayers{background,main,foreground}

% ══════════════════════════════════════════════════════════════
% ── 上半部第 1 行：左（反例一抛物线） + 右（反例二圆周四点） ──
% ══════════════════════════════════════════════════════════════

% ── 1. 左卡片：连续对称抛物线 Y=X^2 ──
\begin{scope}[xshift=0cm, yshift=5.0cm]
  % 卡片底色
  \begin{pgfonlayer}{background}
    \fill[fill=themefg!4!themebg, draw=themegray!30, line width=0.6pt, rounded corners=5pt]
      (-0.2, -0.2) rectangle (7.1, 5.2);
  \end{pgfonlayer}

  % 标题
  \begin{pgfonlayer}{foreground}
    \node[anchor=north west] at (0.0, 4.95) {
      \large \textbf{\color{themecurve}1. 反例一：连续对称 $Y=X^2$}
    };
  \end{pgfonlayer}

  % 坐标系与抛物线
  \begin{pgfonlayer}{background}
    % 坐标轴
    \draw[color=themegray, line width=0.6pt, ->] (0.2, 1.4) -- (2.8, 1.4) node[anchor=west] {\tiny $x$};
    \draw[color=themegray, line width=0.6pt, ->] (1.45, 1.0) -- (1.45, 3.9) node[anchor=south] {\tiny $y$};

    % 抛物线轨迹
    \draw[color=themecurve, line width=1.4pt, domain=-1.15:1.15, samples=40]
      plot ({1.45 + \x}, {1.4 + 1.4*\x*\x});

    % 对称端点
    \fill[color=themeamber] (0.3, 3.25) circle (2.0pt);
    \fill[color=themeamber] (2.6, 3.25) circle (2.0pt);
  \end{pgfonlayer}

  \begin{pgfonlayer}{foreground}
    \node[anchor=north west, color=themecurve] at (0.0, 3.9) {\scriptsize $X \sim U[-1,1]$};
    \node[anchor=north east, color=themeamber] at (2.9, 3.9) {\scriptsize $Y = X^2$};

    % 右侧文字说明
    \node[anchor=north west, text width=3.9cm] at (3.0, 4.3) {
      \small \textbf{协方差为零：}\\[0.25em]
      \scriptsize 由奇对称性：$EX=E[X^3]=0$；\\[0.15em]
      $\operatorname{Cov}(X,Y)=E[X^3]-0 = 0$。\\[0.7em]
      \small \textbf{强确定性非独立：}\\[0.25em]
      \scriptsize $Y=X^2$ 确定依赖，非独立；\\[0.15em]
      $P(Y\le \tfrac{1}{4}\mid X=1)=0 \ne \tfrac{1}{2}$。
    };
  \end{pgfonlayer}
\end{scope}

% ── 2. 右卡片：圆周离散四点 ──
\begin{scope}[xshift=7.7cm, yshift=5.0cm]
  % 卡片底色
  \begin{pgfonlayer}{background}
    \fill[fill=themefg!4!themebg, draw=themegray!30, line width=0.6pt, rounded corners=5pt]
      (-0.2, -0.2) rectangle (7.1, 5.2);
  \end{pgfonlayer}

  % 标题
  \begin{pgfonlayer}{foreground}
    \node[anchor=north west] at (0.0, 4.95) {
      \large \textbf{\color{themeamber}2. 反例二：圆周离散四点}
    };
  \end{pgfonlayer}

  % 几何坐标系与四点
  \begin{pgfonlayer}{background}
    % 虚线辅助圆
    \draw[color=themegray!40, line width=0.6pt, dashed] (1.35, 2.3) circle (0.95cm);

    % 坐标轴
    \draw[color=themegray, line width=0.6pt, ->] (0.2, 2.3) -- (2.6, 2.3) node[anchor=west] {\tiny $x$};
    \draw[color=themegray, line width=0.6pt, ->] (1.35, 1.0) -- (1.35, 3.8) node[anchor=south] {\tiny $y$};

    % 四个离散支撑点 (概率各 1/4)
    \fill[color=themealert] (2.3, 2.3) circle (2.2pt);
    \fill[color=themealert] (1.35, 3.25) circle (2.2pt);
    \fill[color=themealert] (0.4, 2.3) circle (2.2pt);
    \fill[color=themealert] (1.35, 1.35) circle (2.2pt);
  \end{pgfonlayer}

  \begin{pgfonlayer}{foreground}
    \node[anchor=north west, color=themealert] at (2.2, 2.25) {\tiny $(1,0)$};
    \node[anchor=south, color=themealert] at (1.35, 3.35) {\tiny $(0,1)$};
    \node[anchor=north east, color=themealert] at (0.5, 2.25) {\tiny $(-1,0)$};
    \node[anchor=north, color=themealert] at (1.35, 1.25) {\tiny $(0,-1)$};

    % 右侧文字说明
    \node[anchor=north west, text width=4.1cm] at (2.85, 4.3) {
      \small \textbf{乘积恒零与不相关：}\\[0.25em]
      \scriptsize 四个点上均有 $XY \equiv 0$；\\[0.15em]
      $E[XY]=0 \implies \operatorname{Cov}=0$。\\[0.7em]
      \small \textbf{联合事件不独立：}\\[0.25em]
      \scriptsize $P(X=1, Y=1)=0$；\\[0.15em]
      但 $P(X=1)P(Y=1)=\tfrac{1}{16} \ne 0$。
    };
  \end{pgfonlayer}
\end{scope}

% ══════════════════════════════════════════════════════════════
% ── 下半部第 2 行：全宽卡片（独立与不相关辨析及二维正态特权） ──
% ══════════════════════════════════════════════════════════════
\begin{scope}[xshift=0cm, yshift=0.3cm]
  % 卡片底色
  \begin{pgfonlayer}{background}
    \fill[fill=themefg!4!themebg, draw=themegray!30, line width=0.6pt, rounded corners=5pt]
      (-0.2, -0.2) rectangle (14.8, 4.2);
  \end{pgfonlayer}

  % 标题
  \begin{pgfonlayer}{foreground}
    \node[anchor=north west] at (0.0, 3.95) {
      \large \textbf{\color{themealert}3. 概念辨析：独立 vs 不相关与二维正态特权}
    };
  \end{pgfonlayer}

  \begin{pgfonlayer}{foreground}
    % 左侧核心概念对比
    \node[anchor=north west, text width=7.2cm] at (0.1, 3.25) {
      \small $\bullet$ \textbf{\color{themecurve}独立性（强条件）}：$F(x,y) = F_X(x)F_Y(y)$；\\[0.15em]
      \quad 任意可测函数变换均解耦：$E[g(X)h(Y)]=E[g(X)]E[h(Y)]$；\\[0.45em]
      \small $\bullet$ \textbf{\color{themeamber}不相关（弱条件）}：$\operatorname{Cov}(X,Y)=0$；\\[0.15em]
      \quad 仅代表无线性关联，无法排除高阶或对称非线性依赖；\\[0.45em]
      \small $\bullet$ \textbf{\color{themealert}单向蕴含}：独立 $\implies$ 不相关；逆命题不成立。
    };

    % 中间分割线
    \draw[color=themegray!30, line width=0.6pt] (7.4, 0.3) -- (7.4, 3.4);

    % 右侧二维正态与考场警示
    \node[anchor=north west, text width=7.0cm] at (7.7, 3.25) {
      \small \textbf{二维正态分布的等价特权：}\\[0.4em]
      \scriptsize $\bullet$ \textbf{充要等价}：若 $(X,Y)$ \textbf{服从二维联合正态}，则\\[0.2em]
      \quad \normalsize \color{themeamber} 独立 $\iff$ 不相关 $\iff \rho_{XY}=0$\\[0.4em]
      \scriptsize $\bullet$ \textbf{考场核心避坑陷阱}：\\[0.15em]
      \quad \textbf{边缘正态 $\nRightarrow$ 联合正态}；若仅知 $X \sim N, Y \sim N$，\\[0.15em]
      \quad 无法直接由不相关推导独立，必须验证联合正态性！
    };
  \end{pgfonlayer}
\end{scope}

% ══════════════════════════════════════════════════════════════
% ── 底部全图全局总结 (无框轻量排印) ──
% ══════════════════════════════════════════════════════════════
\begin{pgfonlayer}{foreground}
  \node[anchor=north, text=themefg, text width=14.8cm, align=center] at (7.3, -0.2) {
    \footnotesize \textbf{依赖拓扑真理：} 独立性是全分布事件分解（强约束）；不相关 $\operatorname{Cov}=0$ 仅消除了一阶线性趋势；对称几何抵消不能掩盖函数确定依赖；只有在二维联合正态下二者等价。
  };
\end{pgfonlayer}

\end{tikzpicture}
%
}
\end{handbookdiagram}

\subsection{协方差矩阵与线性组合}

对随机向量 $\boldsymbol X=(X_1,\ldots,X_p)^{\mathsf T}$，
\[
\boldsymbol\mu=E(\boldsymbol X),\qquad
\boxed{\boldsymbol\Sigma=E[(\boldsymbol X-\boldsymbol\mu)(\boldsymbol X-\boldsymbol\mu)^{\mathsf T}]}.
\]
第 $(i,j)$ 个元素是 $\operatorname{Cov}(X_i,X_j)$。矩阵对称半正定，因为对任意 $\boldsymbol a$：
\[
\boldsymbol a^{\mathsf T}\boldsymbol\Sigma\boldsymbol a
=\operatorname{Var}(\boldsymbol a^{\mathsf T}\boldsymbol X)\geq0.
\]
这把“方差如何随线性组合变化”压缩成一个矩阵算子，是多维正态、回归和统计推断的接口；本册只拥有其概率定义，统计估计留给 Topic 07/08。

\subsection{线性组合方差的 Rayleigh 商接口}

当题目把线性组合的系数限制为单位长度，例如
\[
a^2+b^2=1,\qquad T=aX+bY,
\]
目标已经不是“展开一次方差”，而是要在所有观察方向中寻找波动最大的方向。令
\[
\boldsymbol a=\begin{pmatrix}a\\b\end{pmatrix},\qquad
\boldsymbol\Sigma=\begin{pmatrix}
\operatorname{Var}(X)&\operatorname{Cov}(X,Y)\\
\operatorname{Cov}(X,Y)&\operatorname{Var}(Y)
\end{pmatrix},
\]
则
\[
\operatorname{Var}(aX+bY)=\boldsymbol a^{\mathsf T}\boldsymbol\Sigma\boldsymbol a,
\qquad \boldsymbol a^{\mathsf T}\boldsymbol a=1.
\]
对称协方差矩阵的 Rayleigh 商满足
\[
\lambda_{\min}(\boldsymbol\Sigma)\leq
\boldsymbol a^{\mathsf T}\boldsymbol\Sigma\boldsymbol a\leq
\lambda_{\max}(\boldsymbol\Sigma),
\]
等号方向分别是对应特征向量。因此“单位系数下方差最大”应先转成最大特征值问题，而不是把一个系数硬设为正后做单变量求导。

\begin{methodbox}{二维标准化正态的直接推导}
若 $\operatorname{Var}(X)=\operatorname{Var}(Y)=1$ 且 $\operatorname{Cov}(X,Y)=\rho$，则
\[
\boldsymbol\Sigma=\begin{pmatrix}1&\rho\\\rho&1\end{pmatrix},\qquad
\lambda_{\pm}=1\pm\rho.
\]
由于最大特征值为 $1+\lvert\rho\rvert$，
\[
\max_{a^2+b^2=1}\operatorname{Var}(aX+bY)=1+\lvert\rho\rvert.
\]
等号方向是 $a=b$ 或 $a=-b$ 的特征方向，具体选择由 $\rho$ 的符号决定。
\end{methodbox}

边界检查：协方差矩阵必须半正定；若题目没有单位长度约束，方差通常没有有限最大值；若系数约束改成一般二次型，应转为广义 Rayleigh 商并检查约束矩阵正定性。该接口连接 Topic 05 的概率摘要与线性代数 Topic 05/06 的谱、正定结构，但不把协方差矩阵重新归入线代 Owner。

\section{第四层：条件期望、迭代与方差分解}

\subsection{条件期望的两种读法}

离散 $Y$ 下，
\[
E(X\mid Y=y)=\sum_xxP(X=x\mid Y=y),
\]
连续情形下，
\[
E(X\mid Y=y)=\int xf_{X\mid Y}(x\mid y)\,dx.
\]
更一般地，$E(X\mid\mathcal G)$ 是对信息 $\mathcal G$ 可测且满足
\[
E[I_AX]=E[I_AE(X\mid\mathcal G)]\qquad(A\in\mathcal G)
\]
的随机变量。直观上，它是在已知信息下对 $X$ 的最佳平方可积预测；代价是条件层之间的信息被保留，而非条件期望只保留总体平均。

\subsection{全期望、全方差与全协方差}

迭代期望（全期望）为
\[
\boxed{E[E(X\mid Y)]=E(X)}.
\]
全方差公式：
\[
\boxed{\operatorname{Var}(X)=E[\operatorname{Var}(X\mid Y)]+\operatorname{Var}[E(X\mid Y)]}.
\]
第一项是层内平均噪声，第二项是层间均值差异；两项均非负，不能漏掉任何一项。二元推广为
\[
\boxed{\operatorname{Cov}(X,Z)=E[\operatorname{Cov}(X,Z\mid Y)]+\operatorname{Cov}(E[X\mid Y],E[Z\mid Y])}.
\]
若条件独立，则第一项为零，但第二项仍可能不为零；这解释了条件独立不等于无条件独立。

\begin{methodbox}{全方差公式的计算推导}
记 $m(Y)=E(X\mid Y)$。把总偏差拆成
\[
X-E(X)=\bigl(X-m(Y)\bigr)+\bigl(m(Y)-E(X)\bigr).
\]
平方后取期望，第一项给 $E[\operatorname{Var}(X\mid Y)]$，第二项给 $\operatorname{Var}(m(Y))$；交叉项为零，因为
\[
E[X-m(Y)\mid Y]=0.
\]
因此全方差不是记忆公式，而是“层内噪声 + 层间均值差异”的正交分解。全协方差完全沿同一拆分，只需把平方换成交叉乘积。
\end{methodbox}

\begin{handbookdiagram}
\resizebox{0.92\linewidth}{!}{%
% ── V11 全方差分解与分层预测（大空间呼吸感 2 行布局） ──
\begin{tikzpicture}[
  scale=1.0,
  >=Stealth,
  font=\small,
  color=themefg,
  text=themefg
]

\pgfdeclarelayer{background}
\pgfdeclarelayer{foreground}
\pgfsetlayers{background,main,foreground}

% ══════════════════════════════════════════════════════════════
% ── 上半部第 1 行：左（正交投影） + 右（全方差分解条带） ──
% ══════════════════════════════════════════════════════════════

% ── 1. 左卡片：条件期望最优正交投影 ──
\begin{scope}[xshift=0cm, yshift=5.0cm]
  % 卡片底色
  \begin{pgfonlayer}{background}
    \fill[fill=themefg!4!themebg, draw=themegray!30, line width=0.6pt, rounded corners=5pt]
      (-0.2, -0.2) rectangle (7.1, 5.2);
  \end{pgfonlayer}

  % 标题
  \begin{pgfonlayer}{foreground}
    \node[anchor=north west] at (0.0, 4.95) {
      \large \textbf{\color{themecurve}1. 条件期望：最优正交投影}
    };
  \end{pgfonlayer}

  % 几何投影图与说明左右分栏
  \begin{pgfonlayer}{background}
    % 子空间基底
    \fill[fill=themegray!15!themebg, draw=themegray!40, line width=0.6pt]
      (0.3, 0.8) -- (3.0, 0.8) -- (2.4, 2.4) -- (0.6, 2.4) -- cycle;

    % 目标随机变量 Y 点
    \fill[color=themealert] (1.65, 3.8) circle (2.2pt);
    % 投影点 E[Y|X]
    \fill[color=themecurve] (1.65, 1.6) circle (2.2pt);

    % 正交投影垂线
    \draw[color=themealert, line width=1.1pt, dashed] (1.65, 3.8) -- (1.65, 1.6);
    % 直角符号
    \draw[color=themegray, line width=0.6pt] (1.65, 1.9) -- (1.9, 1.9) -- (1.9, 1.6);
  \end{pgfonlayer}

  \begin{pgfonlayer}{foreground}
    \node[anchor=west, color=themealert] at (1.75, 3.8) {\normalsize $Y$};
    \node[anchor=north, color=themecurve] at (1.65, 1.45) {\small $E[Y\mid X]$};
    \node[anchor=north east, color=themegray] at (3.0, 2.8) {\scriptsize $\mathcal{L}^2(\sigma(X))$};

    % 右侧文字说明
    \node[anchor=north west, text width=3.7cm] at (3.25, 4.3) {
      \small \textbf{最优均方预测：}\\[0.25em]
      \scriptsize $\arg\min_{g} E\bigl[(Y-g(X))^2\bigr] = E[Y\mid X]$\\[0.7em]
      \small \textbf{残差正交性：}\\[0.25em]
      \scriptsize 对任意可测函数 $h(X)$，有\\[0.2em]
      $E\bigl[(Y-E[Y\mid X])h(X)\bigr] = 0$。
    };
  \end{pgfonlayer}
\end{scope}

% ── 2. 右卡片：全方差组内组间分解 ──
\begin{scope}[xshift=7.7cm, yshift=5.0cm]
  % 卡片底色
  \begin{pgfonlayer}{background}
    \fill[fill=themefg!4!themebg, draw=themegray!30, line width=0.6pt, rounded corners=5pt]
      (-0.2, -0.2) rectangle (7.1, 5.2);
  \end{pgfonlayer}

  % 标题
  \begin{pgfonlayer}{foreground}
    \node[anchor=north west] at (0.0, 4.95) {
      \large \textbf{\color{themeamber}2. 全方差：组间与组内正交分解}
    };
  \end{pgfonlayer}

  % 方差树状分解与说明
  \begin{pgfonlayer}{background}
    \draw[color=themegray, line width=0.8pt, ->] (1.5, 3.0) -- (0.7, 2.1);
    \draw[color=themegray, line width=0.8pt, ->] (1.5, 3.0) -- (2.3, 2.1);
  \end{pgfonlayer}

  \begin{pgfonlayer}{foreground}
    \node[color=themefg] at (1.5, 3.4) {\small \textbf{总方差 $\operatorname{Var}(Y)$}};
    \node[color=themecurve, align=center] at (0.7, 1.45) {\scriptsize \textbf{组间方差}\\[0.2em]\tiny $\operatorname{Var}(E[Y\mid X])$};
    \node[color=themeamber, align=center] at (2.3, 1.45) {\scriptsize \textbf{组内方差}\\[0.2em]\tiny $E[\operatorname{Var}(Y\mid X)]$};

    % 右侧文字说明
    \node[anchor=north west, text width=3.7cm] at (3.25, 4.3) {
      \small \textbf{全方差分解：}\\[0.25em]
      \scriptsize $\operatorname{Var}(Y) = \operatorname{Var}(E[Y|X]) + E[\operatorname{Var}(Y|X)]$\\[0.7em]
      \small \textbf{全期望底座：}\\[0.25em]
      \scriptsize $E\bigl[E[Y\mid X]\bigr] = E[Y]$。
    };
  \end{pgfonlayer}
\end{scope}

% ══════════════════════════════════════════════════════════════
% ── 下半部第 2 行：全宽卡片（实战模型：随机步数和 Wald 恒等式） ──
% ══════════════════════════════════════════════════════════════
\begin{scope}[xshift=0cm, yshift=0.3cm]
  % 卡片底色
  \begin{pgfonlayer}{background}
    \fill[fill=themefg!4!themebg, draw=themegray!30, line width=0.6pt, rounded corners=5pt]
      (-0.2, -0.2) rectangle (14.8, 4.2);
  \end{pgfonlayer}

  % 标题
  \begin{pgfonlayer}{foreground}
    \node[anchor=north west] at (0.0, 3.95) {
      \large \textbf{\color{themealert}3. 实战模型：随机步数和（Wald 恒等式与方差展开）}
    };
  \end{pgfonlayer}

  \begin{pgfonlayer}{foreground}
    % 左侧公式清单（排印宽敞，不拥挤）
    \node[anchor=north west, text width=7.2cm] at (0.1, 3.25) {
      \small $\bullet$ \textbf{\color{themecurve}随机求和模型}：\\[0.15em]
      \quad $S_N = \sum_{i=1}^N X_i$ \quad \scriptsize ($X_i\ \text{i.i.d. 且独立于 } N$)\\[0.45em]
      \small $\bullet$ \textbf{\color{themeamber}Wald 期望}：$E[S_N] = E[N]\cdot E[X]$；\\[0.45em]
      \small $\bullet$ \textbf{\color{themealert}Wald 方差}：\\[0.2em]
      \quad $\operatorname{Var}(S_N) = \underbrace{E[N]\sigma_X^2}_{\text{单步微观波动}} + \underbrace{\mu_X^2 \sigma_N^2}_{\text{步数宏观涨落}}$。
    };

    % 中间分割线
    \draw[color=themegray!30, line width=0.6pt] (7.4, 0.3) -- (7.4, 3.4);

    % 右侧物理机制与警示
    \node[anchor=north west, text width=7.0cm] at (7.7, 3.25) {
      \small \textbf{两层波动机制与考场警示：}\\[0.4em]
      \scriptsize $\bullet$ \textbf{组内微观项} $E[N]\sigma_X^2$：单步随机变量 $X_i$ 自身微观涨落；\\[0.35em]
      \scriptsize $\bullet$ \textbf{组间宏观项} $\mu_X^2 \sigma_N^2$：步数 $N$ 离散跳跃引发的宏观波动；\\[0.35em]
      \scriptsize $\bullet$ \textbf{常数步警示}：即使 $\sigma_X^2=0$，总方差仍为 $c^2\sigma_N^2 \ne 0$！
    };
  \end{pgfonlayer}
\end{scope}

% ══════════════════════════════════════════════════════════════
% ── 底部全图全局总结 (无框轻量排印) ──
% ══════════════════════════════════════════════════════════════
\begin{pgfonlayer}{foreground}
  \node[anchor=north, text=themefg, text width=14.8cm, align=center] at (7.3, -0.2) {
    \footnotesize \textbf{全方差正交哲学：} 条件期望 $E[Y\mid X]$ 是最优均方预测；全方差公式将总波动拆为 \textbf{各层均值的组间方差} 与 \textbf{各层内部波动的期望}；随机步数求和必须计入步数方差。
  };
\end{pgfonlayer}

\end{tikzpicture}
%
}
\end{handbookdiagram}

\begin{methodbox}{分布函数/密度的凸混合：先恢复选择层，再直接算矩}
若
\[
F_X(x)=\sum_{j=1}^m w_jF_j(x),\qquad w_j\ge0,\quad \sum_jw_j=1,
\]
或在有密度时
\[
f_X(x)=\sum_{j=1}^m w_jf_j(x),
\]
就把它解释成：先抽取离散层 $G$，$P(G=j)=w_j$，再令 $X\mid G=j\sim F_j$。因此只问数字特征时，不必先把混合 CDF 求导、再从头积分：若各层相应矩存在，直接有
\[
E(X)=\sum_jw_j\mu_j,
\qquad
E(X^2)=\sum_jw_jE(X_j^2),
\]
\[
\operatorname{Var}(X)=\sum_jw_j\sigma_j^2
+\sum_jw_j(\mu_j-\mu)^2,
\qquad \mu=\sum_jw_j\mu_j.
\]
第一项是层内方差的加权平均，第二项是层间均值波动。旧稿中
$F(x)=0.3\Phi(x)+0.7\Phi((x-1)/2)$ 一类题，第一动作应识别为“以 $0.3$ 选择 $N(0,1)$、以 $0.7$ 选择 $N(1,4)$”，所以期望直接是 $0.3\cdot0+0.7\cdot1$；只有题目继续问完整密度时才需要求导。
\end{methodbox}

\begin{examplebox}{真题对照：随机选择的混合 vs 取平均}
这里沿用全册记号：$X\indep Y$ 表示随机变量概率独立；单符号 $\perp$ 只用于几何正交。
2014 Q八 中，$Y_1$ 以概率 $1/2$ 选择 $X_1$ 或 $X_2$，而 $Y_2=(X_1+X_2)/2$。令 $G\sim\operatorname{Bernoulli}(1/2)$ 独立于 $(X_1,X_2)$，则 $Y_1=GX_1+(1-G)X_2$。条件分解给出
\[
\begin{aligned}
E(Y_1)&=\tfrac12(E X_1+E X_2),\\
\operatorname{Var}(Y_1)&=\tfrac12(\operatorname{Var}X_1+\operatorname{Var}X_2)
 +\tfrac14(E X_1-E X_2)^2,\\
\operatorname{Var}(Y_2)&=\tfrac14(\operatorname{Var}X_1+\operatorname{Var}X_2)
\quad(X_1\indep X_2).
\end{aligned}
\]
额外的层间均值差异使混合的方差通常大于平均值的方差；“均值相同”不意味着波动相同。

2020 Q二十二进一步构造 $Y=X_3X_1+(1-X_3)X_2$，其中 $X_1,X_2$ 为独立标准正态，$X_3$ 等概率取 $0,1$。分层可知 $Y\sim N(0,1)$，但
\[
\operatorname{Cov}(X_1,Y)=E(X_1Y)=\tfrac12E(X_1^2)+\tfrac12E(X_1X_2)=\tfrac12\ne0.
\]
所以即使两个边缘都标准正态，也不能从“看起来同分布”推出独立；必须回到联合生成链。
\end{examplebox}

\begin{examplebox}{真题动作链：2024 Q九 的条件均匀与全协方差}
设 $X$ 的密度为 $f_X(x)=2(1-x)$（$0<x<1$），且给定 $X=x$ 时
$Y\mid X=x\sim U(x,1)$。条件分布的区间随 $x$ 改变，不能把 $Y$
直接当作固定区间上的均匀变量。先取条件均值：
\[
E(Y\mid X=x)=\frac{x+1}{2},\qquad EX=\int_0^1 2x(1-x)\,dx=\frac13,
\]
\[
DX=\int_0^1x^2\,2(1-x)\,dx-\left(\frac13\right)^2=\frac1{18}.
\]
全协方差公式在这里写成
\[
\operatorname{Cov}(X,Y)
=E[\operatorname{Cov}(X,Y\mid X)]
 +\operatorname{Cov}(E[X\mid X],E[Y\mid X]).
\]
第一项为 $0$，因为给定 $X$ 后 $X$ 是常数；第二项为
\[
\operatorname{Cov}\left(X,\frac{X+1}{2}\right)=\frac12DX=\frac1{36}.
\]
若题目改问 $DY$，则不能在此停止，还要使用
\[
DY=E[\operatorname{Var}(Y\mid X)]+\operatorname{Var}\left(\frac{X+1}{2}\right),
\qquad \operatorname{Var}(Y\mid X=x)=\frac{(1-x)^2}{12},
\]
把层内均匀噪声与层间均值波动分别记账。\kw{模型句：}条件均值保留了由 $X$ 传递到 $Y$ 的线性信号；条件方差描述给定 $X$ 后剩余的层内噪声。
\end{examplebox}

\section{第五层：由矩推出概率界}

\subsection{Cauchy--Schwarz 与 Jensen}

若二阶矩存在，
\[
E\lvert XY\rvert\leq\sqrt{E(X^2)E(Y^2)},\qquad \lvert\operatorname{Cov}(X,Y)\rvert\leq\sigma_X\sigma_Y.
\]
若 $\varphi$ 为凸函数且期望存在，Jensen 不等式为
\[
\boxed{\varphi(EX)\leq E[\varphi(X)]}.
\]
严格凸时等号通常要求 $X$ 几乎处处为常数。典型地，$[E(X)]^2\leq E(X^2)$。

\subsection{Markov、Chebyshev 与 Cantelli}

若 $X\geq0,a>0$，Markov：
\[
\boxed{P(X\geq a)\leq\frac{E(X)}a}.
\]
若 $E(X)=\mu,\operatorname{Var}(X)=\sigma^2<\infty$，Chebyshev：
\[
\boxed{P(\lvert X-\mu\rvert\geq\varepsilon)\leq\frac{\sigma^2}{\varepsilon^2}}.
\]
单侧偏离用 Cantelli：
\[
\boxed{P(X-\mu\geq a)\leq\frac{\sigma^2}{\sigma^2+a^2}},\qquad a>0.
\]
选择规则是：非负量先看 Markov；双侧偏差看 Chebyshev；明确单侧方向时检查 Cantelli。它们只给上界，不给事件的精确概率；Chebyshev 连接 Topic 06 的弱大数律，但极限定理证明不在本册。

\begin{methodbox}{概率界的最小证明链}
若 $X\ge0$，事件指标满足 $I_{\{X\ge a\}}\le X/a$，取期望即得 Markov。对一般 $X$，把 $\lvert X-\mu\rvert\ge\varepsilon$ 应用于非负量 $(X-\mu)^2$，得到 Chebyshev。单侧 Cantelli 则对任意 $c>0$ 写出
\[
P(X-\mu\ge a)\le\frac{E[(X-\mu+c)^2]}{(a+c)^2}
=\frac{\sigma^2+c^2}{(a+c)^2},
\]
再对 $c$ 优化，最小值在 $c=\sigma^2/a$ 处给出
\[
P(X-\mu\ge a)\le\frac{\sigma^2}{\sigma^2+a^2}.
\]
三者都只由已知矩推出上界；若题目给出完整分布，不能用界替代精确概率。
\end{methodbox}

\section{常见分布的数字特征速查表}

\begin{longtable}{L{3.5cm}L{4.0cm}L{4.5cm}}
\toprule
\textbf{分布模型} & \textbf{期望 $E(X)$} & \textbf{方差 $\operatorname{Var}(X)$}\\
\midrule
0--1 分布 $B(1,p)$ & $p$ & $p(1-p)$\\
二项分布 $B(n,p)$ & $np$ & $np(1-p)$\\
泊松分布 $P(\lambda)$ & $\lambda$ & $\lambda$\\
几何分布 $G(p)$ & $1/p$ & $(1-p)/p^2$\\
均匀分布 $U(a,b)$ & $(a+b)/2$ & $(b-a)^2/12$\\
指数分布 $\operatorname{Exp}(\lambda)$ & $1/\lambda$ & $1/\lambda^2$\\
正态分布 $N(\mu,\sigma^2)$ & $\mu$ & $\sigma^2$\\
$\chi^2(n)$ 分布 & $n$ & $2n$\\
$t(n)$ 分布 ($n>2$) & 0 & $n/(n-2)$\\
$F(m,n)$ 分布 ($n>2$) & $n/(n-2)$ & 较复杂\\
Gamma $(\alpha,\lambda)$（率参数） & $\alpha/\lambda$ & $\alpha/\lambda^2$\\
Cauchy $(x_0,\gamma)$ & 不存在 & 不存在\\
\bottomrule
\end{longtable}

指数、Gamma、卡方、正态等分布的生成关系与卷积机制归 Topic 04；本表只提供数字特征接口。$t$ 分布的方差要求自由度 $n>2$，$F(m,n)$ 的期望要求 $n>2$，不同参数约定必须先核对。

\section{第六层：总体特征与样本统计接口}

本册讨论总体随机变量的特征；统计学中的样本量化对象由 Topic 07 负责。对应关系如下：
\begin{center}
\begin{tabularx}{\linewidth}{L{3.1cm}L{4.2cm}Y}
\toprule
总体特征 & 常用样本对象 & 归属边界\\
\midrule
$\mu=E(X)$ & $\bar X=n^{-1}\sum_iX_i$ & 其抽样分布与无偏性在 Topic 07\\
$\sigma^2=\operatorname{Var}(X)$ & $S^2=(n-1)^{-1}\sum_i(X_i-\bar X)^2$ & $n-1$ 修正及正态平方和在 Topic 07\\
$E(X^k)$ & $A_k=n^{-1}\sum_iX_i^k$ & 矩估计在 Topic 08 调用\\
$E[(X-\mu)^k]$ & $B_k=n^{-1}\sum_i(X_i-\bar X)^k$ & 样本中心矩不自动无偏\\
$\operatorname{Cov}(X,Y)$ & 样本协方差 & 总体定义在本册，估计性质在 Topic 08\\
$\rho_{XY}$ & 样本相关系数 & 仅描述线性关联，不能替代独立性检验\\
\bottomrule
\end{tabularx}
\end{center}

\section{第七层：统一解题工作流与边界检查}

\begin{enumerate}
  \item 先锁定对象：是总体分布的期望、方差、协方差、高阶矩，还是样本统计量；不要把 $S^2$ 当作总体 $\sigma^2$。
  \item 检查存在性：$E\lvert X\rvert<\infty$ 才能谈有限期望，$E(X^2)<\infty$ 才能使用方差和相关系数；重尾模型先停在边界检查。
  \item 选择最短表示：函数期望优先 LOTUS，计数优先示性变量，分层信息优先全期望/全方差，线性组合优先协方差展开。
  \item 只在需要时使用独立：期望线性性不需要独立；乘积拆分、卷积、交叉协方差消失才需要相应独立或不相关条件。
  \item 做结果校验：方差非负，$\lvert\rho\rvert\leq1$，条件方差分解两项非负，单位与尺度变换一致，概率界不超过 $1$。
\end{enumerate}

\begin{methodbox}{高频结构到动作}
\begin{tabularx}{\linewidth}{L{4.3cm}Y}
\toprule
题面信号 & 首选动作与停止条件\\
\midrule
“满足条件的个数/至少一个” & 写 $N=\sum I_{A_i}$；期望阶段不展开独立性，只有题目问方差时再处理交叉项。\\
“非负寿命、等待时间” & 先试尾积分/尾和；若尾部积分发散，停止使用有限期望公式。\\
“已知某类/某阶段/某组” & 写条件期望，先算层内再按层权重合成；全方差必须保留层内与层间两项。\\
“线性组合的方差” & 展开协方差矩阵；确认独立、两两不相关或仅有一般相关。\\
“只给均值和方差求概率上界” & 选择 Markov/Chebyshev/Cantelli；明确这是界而非精确值。\\
“判断独立或相关” & 零协方差只给线性信息；回到联合分布，联合正态才可用等价结论。\\
\bottomrule
\end{tabularx}
\end{methodbox}

\begin{warnbox}{必须保留的反向边界}
相同数字特征不确定分布；标准化不产生正态性；高阶矩不存在时偏度/峰度无意义；有偏估计量不能仅凭“有偏”判废，需比较 MSE、效率和相合性；总体数字特征与样本统计量的数值形式相似，但随机性与抽样分布不同。
\end{warnbox}

\section{贯穿母例：分层测量中的位置、波动与依赖}

\begin{examplebox}{母例：先分组再测量}
设 $G\sim\operatorname{Bern}(p)$ 表示测量来自普通组（$G=0$）还是高位组（$G=1$）。给定组别后，测量值 $X$ 满足
\[
E(X\mid G=0)=1,\quad E(X\mid G=1)=4,\quad
\operatorname{Var}(X\mid G=0)=\operatorname{Var}(X\mid G=1)=1.
\]
先做条件分解，而不是把两组混成一个分布：
\[
E(X)=E[E(X\mid G)]=1+3p,
\]
\[
\operatorname{Var}(X)=E[\operatorname{Var}(X\mid G)]
 +\operatorname{Var}(E[X\mid G])
 =1+9p(1-p).
\]
若把组别写成示性变量 $G$，则
\[
\operatorname{Cov}(X,G)=\operatorname{Cov}(E[X\mid G],G)=3p(1-p).
\]
当只知道均值和方差而要求偏离概率时，最多先给出 Chebyshev 上界；不能把这个上界当作混合分布的精确概率。这个例子把本册的四个动作串起来：LOTUS/全期望给位置，全方差给层内与层间波动，协方差识别组别依赖，不等式提供信息不足时的保守结论。
\end{examplebox}

\section{一页压缩}

\begin{corebox}{一句话总结}
\[
\boxed{\text{期望是线性加权，方差是波动平方，协方差度量线性交叠，全期望分解层级复杂性。}}
\]
\end{corebox}

\begin{examplebox}{先做逐点恒等式，再做期望：2011 真题}
若 $U=\max\{X,Y\}$、$V=\min\{X,Y\}$，则对每个样本点都有 $UV=XY$。因此
\[
E(UV)=E(XY)=E(X)E(Y)
\]
（最后一步才使用独立性）。只求数字特征时，先寻找逐点恒等式和 LOTUS，通常不需要构造次序统计量的完整联合分布。
\end{examplebox}

\begin{examplebox}{早期真题：目标只问期望时直接走 LOTUS（1992）}
若 $X\sim\operatorname{Exp}(1)$，则
\[
E(X+e^{-2X})=E(X)+E(e^{-2X})=1+\int_0^\infty e^{-3x}dx=\frac43.
\]
题目只要求数字特征时，不必先构造 $X+e^{-2X}$ 的完整分布；这是“输出层决定计算成本”的直接证据。
\end{examplebox}

\end{document}
